\markdownRendererDocumentBegin
\markdownRendererSectionBegin
\markdownRendererHeadingOne{Programación Dinamica}\markdownRendererInterblockSeparator
{}\markdownRendererSectionBegin
\markdownRendererSectionBegin
\markdownRendererSectionBegin
\markdownRendererHeadingFour{Identificar el subproblema}\markdownRendererInterblockSeparator
{}\markdownRendererUlBeginTight
\markdownRendererUlItem Encuentra la \markdownRendererStrongEmphasis{pequeña versión del problema} que puedes resolver y que se reutiliza.\markdownRendererUlItemEnd 
\markdownRendererUlItem Pregúntate: "Si conozco la solución a los casos más pequeños, ¿puedo construir la solución al caso grande?"\markdownRendererUlItemEnd 
\markdownRendererUlItem Ejemplo: En Coin Change, subproblema = "mínimo número de monedas para sumar x".\markdownRendererUlItemEnd 
\markdownRendererUlEndTight \markdownRendererParagraphSeparator
{}
\markdownRendererSectionEnd \markdownRendererSectionBegin
\markdownRendererHeadingFour{Formular la recurrencia}\markdownRendererInterblockSeparator
{}\markdownRendererUlBeginTight
\markdownRendererUlItem Escribe cómo se relaciona un caso con los casos más pequeños.\markdownRendererUlItemEnd 
\markdownRendererUlItem Generalmente tiene la forma:\markdownRendererSoftLineBreak
{}\markdownRendererDollarSign{}\markdownRendererDollarSign{}\markdownRendererSoftLineBreak
{}\markdownRendererBackslash{}text\markdownRendererLeftBrace{}dp\markdownRendererRightBrace{}[x] = \markdownRendererBackslash{}text\markdownRendererLeftBrace{}función de dp[casos más pequeños]\markdownRendererRightBrace{}\markdownRendererSoftLineBreak
{}\markdownRendererDollarSign{}\markdownRendererDollarSign{}\markdownRendererUlItemEnd 
\markdownRendererUlItem Ejemplo:\markdownRendererSoftLineBreak
{}\markdownRendererDollarSign{}\markdownRendererDollarSign{}\markdownRendererSoftLineBreak
{}\markdownRendererBackslash{}text\markdownRendererLeftBrace{}solve\markdownRendererRightBrace{}(x) = \markdownRendererBackslash{}min\markdownRendererUnderscore{}\markdownRendererLeftBrace{}c \markdownRendererBackslash{}in coins\markdownRendererRightBrace{} (\markdownRendererBackslash{}text\markdownRendererLeftBrace{}solve\markdownRendererRightBrace{}(x-c)+1)\markdownRendererSoftLineBreak
{}\markdownRendererDollarSign{}\markdownRendererDollarSign{}\markdownRendererUlItemEnd 
\markdownRendererUlEndTight \markdownRendererParagraphSeparator
{}
\markdownRendererSectionEnd \markdownRendererSectionBegin
\markdownRendererHeadingFour{Definir casos base}\markdownRendererInterblockSeparator
{}\markdownRendererUlBeginTight
\markdownRendererUlItem Piensa en \markdownRendererStrongEmphasis{valores triviales} que no dependen de otros subproblemas.\markdownRendererUlItemEnd 
\markdownRendererUlItem Ejemplo Coin Change:\markdownRendererInterblockSeparator
{}\markdownRendererUlBeginTight
\markdownRendererUlItem \markdownRendererCodeSpan{solve(0) = 0} → no necesitas monedas para sumar 0\markdownRendererUlItemEnd 
\markdownRendererUlItem \markdownRendererCodeSpan{solve(x<0) = ∞} → imposible\markdownRendererUlItemEnd 
\markdownRendererUlEndTight \markdownRendererUlItemEnd 
\markdownRendererUlEndTight \markdownRendererInterblockSeparator
{}
\markdownRendererSectionEnd \markdownRendererSectionBegin
\markdownRendererHeadingFour{Memoización o tabla iterativa}\markdownRendererInterblockSeparator
{}\markdownRendererUlBeginTight
\markdownRendererUlItem \markdownRendererStrongEmphasis{Memoización:} Guardar resultados de subproblemas recursivos en un array.\markdownRendererUlItemEnd 
\markdownRendererUlItem \markdownRendererStrongEmphasis{Iterativa (bottom-up):} Llenar un array desde los casos base hasta el objetivo.\markdownRendererUlItemEnd 
\markdownRendererUlItem Ventaja iterativa: menos sobrecarga de llamadas recursivas y más control sobre complejidad.\markdownRendererUlItemEnd 
\markdownRendererUlEndTight \markdownRendererInterblockSeparator
{}
\markdownRendererSectionEnd \markdownRendererSectionBegin
\markdownRendererHeadingFour{Reconstrucción de la solución}\markdownRendererInterblockSeparator
{}\markdownRendererUlBeginTight
\markdownRendererUlItem A veces necesitamos \markdownRendererStrongEmphasis{no solo el valor óptimo}, sino \markdownRendererStrongEmphasis{cómo construirlo}.\markdownRendererUlItemEnd 
\markdownRendererUlItem Guardar información adicional en arrays auxiliares, por ejemplo:\markdownRendererInterblockSeparator
{}\markdownRendererUlBeginTight
\markdownRendererUlItem \markdownRendererCodeSpan{first[x]} → primer paso para llegar a x\markdownRendererUlItemEnd 
\markdownRendererUlEndTight \markdownRendererUlItemEnd 
\markdownRendererUlItem Luego se recorre hacia atrás para reconstruir la solución.\markdownRendererUlItemEnd 
\markdownRendererUlEndTight \markdownRendererInterblockSeparator
{}
\markdownRendererSectionEnd \markdownRendererSectionBegin
\markdownRendererHeadingFour{Contar soluciones}\markdownRendererInterblockSeparator
{}\markdownRendererUlBeginTight
\markdownRendererUlItem A veces queremos \markdownRendererStrongEmphasis{todas las formas posibles} de alcanzar un estado.\markdownRendererUlItemEnd 
\markdownRendererUlItem La recurrencia cambia de "min" o "max" a "suma":\markdownRendererSoftLineBreak
{}\markdownRendererDollarSign{}\markdownRendererDollarSign{}\markdownRendererSoftLineBreak
{}\markdownRendererBackslash{}text\markdownRendererLeftBrace{}count\markdownRendererRightBrace{}[x] = \markdownRendererBackslash{}sum\markdownRendererUnderscore{}\markdownRendererLeftBrace{}c \markdownRendererBackslash{}in coins\markdownRendererRightBrace{} \markdownRendererBackslash{}text\markdownRendererLeftBrace{}count\markdownRendererRightBrace{}[x-c]\markdownRendererSoftLineBreak
{}\markdownRendererDollarSign{}\markdownRendererDollarSign{}\markdownRendererUlItemEnd 
\markdownRendererUlEndTight \markdownRendererInterblockSeparator
{}
\markdownRendererSectionEnd \markdownRendererSectionBegin
\markdownRendererHeadingFour{Optimización de espacio}\markdownRendererInterblockSeparator
{}\markdownRendererUlBeginTight
\markdownRendererUlItem Muchos DP pueden reducir el tamaño del array si solo dependen de los últimos k valores.\markdownRendererUlItemEnd 
\markdownRendererUlItem Ejemplo: Coin Change solo depende de \markdownRendererCodeSpan{dp[x-c]} → se puede usar un vector de tamaño n+1.\markdownRendererUlItemEnd 
\markdownRendererUlEndTight \markdownRendererInterblockSeparator
{}
\markdownRendererSectionEnd 
\markdownRendererSectionEnd \markdownRendererSectionBegin
\markdownRendererHeadingThree{Patrones Tipicos}\markdownRendererInterblockSeparator
{}
\markdownRendererSectionEnd \markdownRendererSectionBegin
\markdownRendererHeadingThree{Lineal (1D DP)}\markdownRendererInterblockSeparator
{}Subproblema depende de un solo parámetro, ej. suma o longitud.\markdownRendererSoftLineBreak
{}Ejemplo: Coin Change (mínimo número de monedas para sumar x)\markdownRendererParagraphSeparator
{}Recursivo:\markdownRendererSoftLineBreak
{}\markdownRendererDollarSign{}\markdownRendererDollarSign{}\markdownRendererSoftLineBreak
{}\markdownRendererBackslash{}text\markdownRendererLeftBrace{}solve\markdownRendererRightBrace{}(x) =\markdownRendererSoftLineBreak
{}\markdownRendererBackslash{}begin\markdownRendererLeftBrace{}cases\markdownRendererRightBrace{}\markdownRendererSoftLineBreak
{}\markdownRendererBackslash{}infty \markdownRendererAmpersand{} x < 0 \markdownRendererBackslash{}\markdownRendererBackslash{}\markdownRendererSoftLineBreak
{}0 \markdownRendererAmpersand{} x = 0 \markdownRendererBackslash{}\markdownRendererBackslash{}\markdownRendererSoftLineBreak
{}\markdownRendererBackslash{}min\markdownRendererUnderscore{}\markdownRendererLeftBrace{}c \markdownRendererBackslash{}in coins\markdownRendererRightBrace{} (\markdownRendererBackslash{}text\markdownRendererLeftBrace{}solve\markdownRendererRightBrace{}(x-c)+1) \markdownRendererAmpersand{} x > 0\markdownRendererSoftLineBreak
{}\markdownRendererBackslash{}end\markdownRendererLeftBrace{}cases\markdownRendererRightBrace{}\markdownRendererSoftLineBreak
{}\markdownRendererDollarSign{}\markdownRendererDollarSign{}\markdownRendererInterblockSeparator
{}\markdownRendererInputFencedCode{./_markdown_main/b28c9b71dbdc1220a00f7a66d9146df1.verbatim}{cpp}{cpp}\markdownRendererInterblockSeparator
{}
\markdownRendererSectionEnd 
\markdownRendererSectionEnd \markdownRendererSectionBegin
\markdownRendererHeadingTwo{Conjuntos (Bitmask DP)}\markdownRendererInterblockSeparator
{}\markdownRendererStrongEmphasis{Ejemplo:} Traveling Salesman Problem (TSP) – recorrer todos los nodos con mínimo costo.\markdownRendererSoftLineBreak
{}\markdownRendererStrongEmphasis{Idea:} Representamos cada subconjunto con un bitmask; \markdownRendererCodeSpan{dp[mask][i]} = mínimo costo para llegar a nodo \markdownRendererCodeSpan{i} visitando los nodos de \markdownRendererCodeSpan{mask}.\markdownRendererParagraphSeparator
{}Recursivo:\markdownRendererSoftLineBreak
{}\markdownRendererDollarSign{}\markdownRendererDollarSign{}\markdownRendererSoftLineBreak
{}\markdownRendererBackslash{}text\markdownRendererLeftBrace{}dp\markdownRendererRightBrace{}\markdownRendererWarning{Undefined link reference "i"}{Undefined link reference "i"}{Look for the text `[...][i]`.}{Look for the text `[...][i]`.}[mask][i] = \markdownRendererBackslash{}min\markdownRendererUnderscore{}\markdownRendererLeftBrace{}j \markdownRendererBackslash{}in mask, j\markdownRendererBackslash{}ne i\markdownRendererRightBrace{} (\markdownRendererBackslash{}text\markdownRendererLeftBrace{}dp\markdownRendererRightBrace{}\markdownRendererWarning{Undefined link reference "j"}{Undefined link reference "j"}{Look for the text `[...][j]`.}{Look for the text `[...][j]`.}[mask\markdownRendererBackslash{}setminus\markdownRendererLeftBrace{}i\markdownRendererRightBrace{}][j] + cost\markdownRendererWarning{Undefined link reference "i"}{Undefined link reference "i"}{Look for the text `[...][i]`.}{Look for the text `[...][i]`.}[j][i])\markdownRendererSoftLineBreak
{}\markdownRendererDollarSign{}\markdownRendererDollarSign{}\markdownRendererInterblockSeparator
{}\markdownRendererInputFencedCode{./_markdown_main/90148e91c3b1dae36f2b131a2edefa45.verbatim}{cpp}{cpp}\markdownRendererInterblockSeparator
{}\markdownRendererDollarSign{}\markdownRendererBackslash{}clearpage\markdownRendererDollarSign{}\markdownRendererInterblockSeparator
{}
\markdownRendererSectionEnd \markdownRendererSectionBegin
\markdownRendererHeadingTwo{Intervalos}\markdownRendererInterblockSeparator
{}Subproblema depende de un rango \markdownRendererCodeSpan{[i,j]}.\markdownRendererSoftLineBreak
{}\markdownRendererStrongEmphasis{Ejemplo:} Matrix Chain Multiplication – minimizar costo de multiplicación de matrices.\markdownRendererSoftLineBreak
{}\markdownRendererStrongEmphasis{Idea:} Para cada intervalo \markdownRendererCodeSpan{[i,j]}, probamos todos los puntos \markdownRendererCodeSpan{k} donde podemos dividir el problema en dos subintervalos.\markdownRendererParagraphSeparator
{}Recursivo:\markdownRendererSoftLineBreak
{}\markdownRendererDollarSign{}\markdownRendererDollarSign{}\markdownRendererSoftLineBreak
{}\markdownRendererBackslash{}text\markdownRendererLeftBrace{}dp\markdownRendererRightBrace{}\markdownRendererWarning{Undefined link reference "j"}{Undefined link reference "j"}{Look for the text `[...][j]`.}{Look for the text `[...][j]`.}[i][j] =\markdownRendererSoftLineBreak
{}\markdownRendererBackslash{}begin\markdownRendererLeftBrace{}cases\markdownRendererRightBrace{}\markdownRendererSoftLineBreak
{}0 \markdownRendererAmpersand{} i = j \markdownRendererBackslash{}\markdownRendererBackslash{}\markdownRendererSoftLineBreak
{}\markdownRendererBackslash{}min\markdownRendererUnderscore{}\markdownRendererLeftBrace{}k=i\markdownRendererRightBrace{}\markdownRendererCircumflex{}\markdownRendererLeftBrace{}j-1\markdownRendererRightBrace{} (\markdownRendererBackslash{}text\markdownRendererLeftBrace{}dp\markdownRendererRightBrace{}\markdownRendererWarning{Undefined link reference "k"}{Undefined link reference "k"}{Look for the text `[...][k]`.}{Look for the text `[...][k]`.}[i][k] + \markdownRendererBackslash{}text\markdownRendererLeftBrace{}dp\markdownRendererRightBrace{}\markdownRendererWarning{Undefined link reference "j"}{Undefined link reference "j"}{Look for the text `[...][j]`.}{Look for the text `[...][j]`.}[k+1][j] + cost\markdownRendererWarning{Undefined link reference "k"}{Undefined link reference "k"}{Look for the text `[...][k]`.}{Look for the text `[...][k]`.}[i]\markdownRendererWarning{Undefined link reference "j"}{Undefined link reference "j"}{Look for the text `[...][j]`.}{Look for the text `[...][j]`.}[k][j]) \markdownRendererAmpersand{} i < j\markdownRendererSoftLineBreak
{}\markdownRendererBackslash{}end\markdownRendererLeftBrace{}cases\markdownRendererRightBrace{}\markdownRendererSoftLineBreak
{}\markdownRendererDollarSign{}\markdownRendererDollarSign{}\markdownRendererInterblockSeparator
{}\markdownRendererInputFencedCode{./_markdown_main/7a67ed37d11796b1f2b51b2d31b3b16c.verbatim}{cpp}{cpp}\markdownRendererInterblockSeparator
{}
\markdownRendererSectionEnd \markdownRendererSectionBegin
\markdownRendererHeadingTwo{Longest Increasing Subsequence (LIS)}\markdownRendererInterblockSeparator
{}\markdownRendererUlBeginTight
\markdownRendererUlItem Encontrar la subsecuencia creciente más larga en un arreglo.\markdownRendererUlItemEnd 
\markdownRendererUlItem Subproblema: \markdownRendererCodeSpan{length(k)} = largo de la LIS que termina en \markdownRendererCodeSpan{k}.\markdownRendererUlItemEnd 
\markdownRendererUlEndTight \markdownRendererInterblockSeparator
{}Recursivo:\markdownRendererSoftLineBreak
{}\markdownRendererDollarSign{}\markdownRendererDollarSign{}\markdownRendererSoftLineBreak
{}length(k) =\markdownRendererSoftLineBreak
{}\markdownRendererBackslash{}begin\markdownRendererLeftBrace{}cases\markdownRendererRightBrace{}\markdownRendererSoftLineBreak
{}1 \markdownRendererAmpersand{} \markdownRendererBackslash{}text\markdownRendererLeftBrace{}si no existe \markdownRendererRightBrace{} i<k \markdownRendererBackslash{}text\markdownRendererLeftBrace{} con \markdownRendererRightBrace{} array[i]<array[k] \markdownRendererBackslash{}\markdownRendererBackslash{}\markdownRendererSoftLineBreak
{}\markdownRendererBackslash{}max\markdownRendererUnderscore{}\markdownRendererLeftBrace{}i<k, array[i]<array[k]\markdownRendererRightBrace{} (length(i)+1) \markdownRendererAmpersand{} \markdownRendererBackslash{}text\markdownRendererLeftBrace{}en otro caso\markdownRendererRightBrace{}\markdownRendererSoftLineBreak
{}\markdownRendererBackslash{}end\markdownRendererLeftBrace{}cases\markdownRendererRightBrace{}\markdownRendererSoftLineBreak
{}\markdownRendererDollarSign{}\markdownRendererDollarSign{}\markdownRendererInterblockSeparator
{}\markdownRendererInputFencedCode{./_markdown_main/10285fabe9c94a9e96408cb62cf0c020.verbatim}{cpp}{cpp}\markdownRendererInterblockSeparator
{}\markdownRendererDollarSign{}\markdownRendererBackslash{}clearpage\markdownRendererDollarSign{}\markdownRendererInterblockSeparator
{}
\markdownRendererSectionEnd \markdownRendererSectionBegin
\markdownRendererHeadingTwo{Paths in a Grid}\markdownRendererInterblockSeparator
{}\markdownRendererUlBeginTight
\markdownRendererUlItem Encontrar la ruta de suma máxima desde la esquina superior izquierda a la inferior derecha, solo movimientos hacia abajo o derecha.\markdownRendererUlItemEnd 
\markdownRendererUlItem Subproblema: \markdownRendererCodeSpan{sum(y,x)} = suma máxima hasta la celda \markdownRendererCodeSpan{(y,x)}.\markdownRendererUlItemEnd 
\markdownRendererUlEndTight \markdownRendererInterblockSeparator
{}Recursivo:\markdownRendererSoftLineBreak
{}\markdownRendererDollarSign{}\markdownRendererDollarSign{}\markdownRendererSoftLineBreak
{}sum(y,x) = max(sum(y-1,x), sum(y,x-1)) + value\markdownRendererWarning{Undefined link reference "x"}{Undefined link reference "x"}{Look for the text `[...][x]`.}{Look for the text `[...][x]`.}[y][x]\markdownRendererSoftLineBreak
{}\markdownRendererDollarSign{}\markdownRendererDollarSign{}\markdownRendererInterblockSeparator
{}\markdownRendererInputFencedCode{./_markdown_main/351fef0456b4f071016b82fcc116775b.verbatim}{cpp}{cpp}\markdownRendererInterblockSeparator
{}
\markdownRendererSectionEnd \markdownRendererSectionBegin
\markdownRendererHeadingTwo{Knapsack}\markdownRendererInterblockSeparator
{}\markdownRendererUlBeginTight
\markdownRendererUlItem Problema: Dado un conjunto de objetos con pesos \markdownRendererDollarSign{}w[i]\markdownRendererDollarSign{} y valores \markdownRendererDollarSign{}v[i]\markdownRendererDollarSign{}, encontrar el subconjunto de máximo valor sin superar capacidad \markdownRendererDollarSign{}c\markdownRendererDollarSign{}.\markdownRendererUlItemEnd 
\markdownRendererUlItem Subproblema: \markdownRendererDollarSign{}dp\markdownRendererWarning{Undefined link reference "x"}{Undefined link reference "x"}{Look for the text `[...][x]`.}{Look for the text `[...][x]`.}[i][x]\markdownRendererDollarSign{} = máximo valor usando los primeros \markdownRendererDollarSign{}i\markdownRendererDollarSign{} objetos con capacidad \markdownRendererDollarSign{}x\markdownRendererDollarSign{}.\markdownRendererUlItemEnd 
\markdownRendererUlEndTight \markdownRendererInterblockSeparator
{}Recursivo:\markdownRendererSoftLineBreak
{}\markdownRendererDollarSign{}\markdownRendererDollarSign{}\markdownRendererSoftLineBreak
{}dp\markdownRendererWarning{Undefined link reference "x"}{Undefined link reference "x"}{Look for the text `[...][x]`.}{Look for the text `[...][x]`.}[i][x] =\markdownRendererSoftLineBreak
{}\markdownRendererBackslash{}begin\markdownRendererLeftBrace{}cases\markdownRendererRightBrace{}\markdownRendererSoftLineBreak
{}dp\markdownRendererWarning{Undefined link reference "x"}{Undefined link reference "x"}{Look for the text `[...][x]`.}{Look for the text `[...][x]`.}[i-1][x] \markdownRendererAmpersand{} \markdownRendererBackslash{}text\markdownRendererLeftBrace{}si no usamos el objeto i\markdownRendererRightBrace{} \markdownRendererBackslash{}\markdownRendererBackslash{}\markdownRendererSoftLineBreak
{}dp[i-1][x-w[i]] + v[i] \markdownRendererAmpersand{} \markdownRendererBackslash{}text\markdownRendererLeftBrace{}si usamos el objeto i (x \markdownRendererBackslash{}ge w[i])\markdownRendererRightBrace{}\markdownRendererSoftLineBreak
{}\markdownRendererBackslash{}end\markdownRendererLeftBrace{}cases\markdownRendererRightBrace{}\markdownRendererSoftLineBreak
{}\markdownRendererDollarSign{}\markdownRendererDollarSign{}\markdownRendererInterblockSeparator
{}\markdownRendererInputFencedCode{./_markdown_main/3b9d34d9574dfd35a440a835e97a7e24.verbatim}{cpp}{cpp}\markdownRendererInterblockSeparator
{}\markdownRendererSectionBegin
\markdownRendererHeadingThree{Reconstruir Knapsack}\markdownRendererInterblockSeparator
{}\markdownRendererInputFencedCode{./_markdown_main/ec7e5a30a4b4ad83e853c1aad284ca54.verbatim}{cpp}{cpp}\markdownRendererInterblockSeparator
{}\markdownRendererStrongEmphasis{Ejemplo}\markdownRendererInterblockSeparator
{}\markdownRendererInputFencedCode{./_markdown_main/5a4418949961f28a95ec871bb87f96d9.verbatim}{cpp}{cpp}\markdownRendererInterblockSeparator
{}\markdownRendererDollarSign{}\markdownRendererBackslash{}clearpage\markdownRendererDollarSign{}\markdownRendererInterblockSeparator
{}
\markdownRendererSectionEnd 
\markdownRendererSectionEnd \markdownRendererSectionBegin
\markdownRendererHeadingTwo{Elevator / Subset DP}\markdownRendererInterblockSeparator
{}\markdownRendererUlBeginTight
\markdownRendererUlItem Mínimo número de viajes para subir personas con límite de peso.\markdownRendererUlItemEnd 
\markdownRendererUlItem Subproblema: Para cada subconjunto \markdownRendererCodeSpan{S}, calcular \markdownRendererCodeSpan{(rides(S), last(S))}.\markdownRendererUlItemEnd 
\markdownRendererUlEndTight \markdownRendererInterblockSeparator
{}Código:\markdownRendererInterblockSeparator
{}\markdownRendererInputFencedCode{./_markdown_main/2b4ee08eda59b49813416f670dc82c61.verbatim}{cpp}{cpp}\markdownRendererParagraphSeparator
{}\markdownRendererDollarSign{}\markdownRendererBackslash{}clearpage\markdownRendererDollarSign{}\markdownRendererInterblockSeparator
{}
\markdownRendererSectionEnd \markdownRendererSectionBegin
\markdownRendererHeadingTwo{Counting Tilings (Grid DP)}\markdownRendererInterblockSeparator
{}\markdownRendererUlBeginTight
\markdownRendererUlItem Contar formas de llenar un grid \markdownRendererCodeSpan{n x m} con piezas \markdownRendererCodeSpan{1x2} y \markdownRendererCodeSpan{2x1}.\markdownRendererUlItemEnd 
\markdownRendererUlItem Subproblema: \markdownRendererCodeSpan{count(k,x)} = formas de llenar filas \markdownRendererCodeSpan{1..k} con estado \markdownRendererCodeSpan{x} en fila \markdownRendererCodeSpan{k}.\markdownRendererUlItemEnd 
\markdownRendererUlEndTight \markdownRendererInterblockSeparator
{}Idea:\markdownRendererInterblockSeparator
{}\markdownRendererUlBeginTight
\markdownRendererUlItem Representar cada fila como un estado.\markdownRendererUlItemEnd 
\markdownRendererUlItem \markdownRendererCodeSpan{count(k,x)} depende solo de \markdownRendererCodeSpan{count(k-1,y)} donde \markdownRendererCodeSpan{y} es compatible con \markdownRendererCodeSpan{x}.\markdownRendererUlItemEnd 
\markdownRendererUlItem Optimización: representar solo columnas con parte superior de los tiles verticales → 2\markdownRendererCircumflex{}m estados.\markdownRendererUlItemEnd 
\markdownRendererUlEndTight \markdownRendererInterblockSeparator
{}Código:\markdownRendererInterblockSeparator
{}\markdownRendererInputFencedCode{./_markdown_main/093dea29c8e57271575f37723cd08162.verbatim}{cpp}{cpp}
\markdownRendererSectionEnd 
\markdownRendererSectionEnd \markdownRendererDocumentEnd